% GENERATED FILE. Edit canonical lessons, then run npm run generate:books.
% canonical-source-hash: fnv1a64:4f1b7b6a98b5c951
% canonical-lessons: FA-C449-masiryab, FA-C449-donbal-kardan, FA-C449-nazdik-shodan, FA-C449-gozashtan, FA-C449-kuche-pas-kuche

\chapter{Following a route}
\label{ch:fa-following-a-route}
\hlchaptermodality{449}

\begin{chapteropening}
\textbf{By the end of this chapter:} \emph{I can follow a route: the satnav, following, getting closer, passing by and the back streets.}

Complete the last lesson of chapter 449: I can follow a route: the satnav, following, getting closer, passing by and the back streets.
\end{chapteropening}

\section[\texorpdfstring{\fa{مسیریاب} (masiryâb)}{مسیریاب (masiryâb)}]{\fa{مسیریاب} (masiryâb) — a satnav, a GPS; a route-finding app}

\label{lesson:FA-C449-masiryab}



\begin{quote}
Before the new one: say the Persian for an underpass, then the Persian for a flyover.
\end{quote}

\subsection*{You'll want to know: \fa{مسیریاب}}
\textbf{\fa{مسیریاب}} (\emph{masiryâb}) — \textquotedblleft{}a satnav, a GPS; a route-finding app\textquotedblright{}.

\begin{quote}
\fa{بدون} \fa{مسیریاب} \fa{در} \fa{این} \fa{شهر} \fa{گم} \fa{میشویم}.
\end{quote}

\begin{quote}
\emph{Bedun-e masiryâb dar in shahr gom mishavim.}
\end{quote}

\begin{quote}
Without a satnav we get lost in this town.
\end{quote}

\emph{masir} 'route' + \emph{yâb} 'finding', from \emph{yâftan}.

\begin{cousinweb}
\emph{masiryâb} is a modern compound of \emph{masir} 'route', a loan usually traced to Arabic \emph{masīr}, and native \emph{yâb}, the present stem of \emph{yâftan} 'to find'.
\end{cousinweb}

\begin{grammarlens}[title={What you've built}]
One more word for telling, explaining and describing at B1.
\end{grammarlens}

\subsection*{Your turn}
\begin{itemize}
\raggedright
  \item \emph{Say it:} \emph{masiryâb}
  \item \emph{Say it:} \emph{masiryâb}, once more
  \item \emph{Say it:} a sentence of your own with \emph{masiryâb}
  \item \emph{Recall:} say the Persian for an underpass, then the Persian for a flyover, then say \emph{masiryâb} again
\end{itemize}

\begin{tcolorbox}[breakable,colback=teal!4,colframe=teal!35!black,title={Before you move on}]
What does \fa{مسیریاب} mean? (\textquotedblleft{}A satnav, a GPS; a route-finding app\textquotedblright{}.) Say it once more.
\end{tcolorbox}
\section[\texorpdfstring{\fa{دنبال} \fa{کردن} (donbâl kardan)}{دنبال کردن (donbâl kardan)}]{\fa{دنبال} \fa{کردن} (donbâl kardan) — to follow}

\label{lesson:FA-C449-donbal-kardan}



\begin{quote}
Before the new one: say the Persian for a flyover, then the Persian for a satnav.
\end{quote}

\subsection*{You'll want to know: \fa{دنبال} \fa{کردن}}
\textbf{\fa{دنبال} \fa{کردن}} (\emph{donbâl kardan}) — \textquotedblleft{}to follow\textquotedblright{}.

\begin{quote}
\fa{ما} \fa{ماشین} \fa{دوستمان} \fa{را} \fa{تا} \fa{روستا} \fa{دنبال} \fa{کردیم}.
\end{quote}

\begin{quote}
\emph{Mâ mâshin-e dustemân râ tâ rustâ donbâl kardim.}
\end{quote}

\begin{quote}
We followed our friend's car to the village.
\end{quote}

\emph{donbâl} 'behind, the tail' + \emph{kardan}. Also to look for: \emph{donbâl-e kelid migardam}.

\begin{cousinweb}
\emph{donbâl kardan} is built from native \emph{donbâl} 'tail, the back', whose older history is not settled, and \emph{kardan} 'to do', inherited from Middle Persian \emph{kardan} (EtymDB).
\end{cousinweb}

\begin{grammarlens}[title={What you've built}]
One more word for telling, explaining and describing at B1.
\end{grammarlens}

\subsection*{Your turn}
\begin{itemize}
\raggedright
  \item \emph{Say it:} \emph{donbâl kardan}
  \item \emph{Say it:} \emph{donbâl kardan}, once more
  \item \emph{Say it:} a sentence of your own with \emph{donbâl kardan}
  \item \emph{Recall:} say the Persian for a flyover, then the Persian for a satnav, then say \emph{donbâl kardan} again
\end{itemize}

\begin{tcolorbox}[breakable,colback=teal!4,colframe=teal!35!black,title={Before you move on}]
What does \fa{دنبال} \fa{کردن} mean? (\textquotedblleft{}To follow\textquotedblright{}.) Say it once more.
\end{tcolorbox}
\section[\texorpdfstring{\fa{نزدیک} \fa{شدن} (nazdik shodan)}{نزدیک شدن (nazdik shodan)}]{\fa{نزدیک} \fa{شدن} (nazdik shodan) — to approach, to get near}

\label{lesson:FA-C449-nazdik-shodan}



\begin{quote}
Before the new one: say the Persian for a satnav, then the Persian for to follow.
\end{quote}

\subsection*{You'll want to know: \fa{نزدیک} \fa{شدن}}
\textbf{\fa{نزدیک} \fa{شدن}} (\emph{nazdik shodan}) — \textquotedblleft{}to approach, to get near\textquotedblright{}.

\begin{quote}
\fa{قطار} \fa{به} \fa{ایستگاه} \fa{نزدیک} \fa{شد}.
\end{quote}

\begin{quote}
\emph{Qatâr be istgâh nazdik shod.}
\end{quote}

\begin{quote}
The train approached the station.
\end{quote}

\emph{nazdik} 'near' + \emph{shodan}; the place takes \emph{be}.

\begin{cousinweb}
\emph{nazdik shodan} is built from native \emph{nazdik} 'near', usually traced to Middle Persian \emph{nazdīk}, and the native verb \emph{shodan} 'to become'.
\end{cousinweb}

\begin{grammarlens}[title={What you've built}]
One more word for telling, explaining and describing at B1.
\end{grammarlens}

\subsection*{Your turn}
\begin{itemize}
\raggedright
  \item \emph{Say it:} \emph{nazdik shodan}
  \item \emph{Say it:} \emph{nazdik shodan}, once more
  \item \emph{Say it:} a sentence of your own with \emph{nazdik shodan}
  \item \emph{Recall:} say the Persian for a satnav, then the Persian for to follow, then say \emph{nazdik shodan} again
\end{itemize}

\begin{tcolorbox}[breakable,colback=teal!4,colframe=teal!35!black,title={Before you move on}]
What does \fa{نزدیک} \fa{شدن} mean? (\textquotedblleft{}To approach, to get near\textquotedblright{}.) Say it once more.
\end{tcolorbox}
\section[\texorpdfstring{\fa{گذشتن} (gozashtan)}{گذشتن (gozashtan)}]{\fa{گذشتن} (gozashtan) — to pass (by, through); to go by (of time)}

\label{lesson:FA-C449-gozashtan}



\begin{quote}
Before the new one: say the Persian for to follow, then the Persian for to approach.
\end{quote}

\subsection*{You'll want to know: \fa{گذشتن}}
\textbf{\fa{گذشتن}} (\emph{gozashtan}) — \textquotedblleft{}to pass (by, through); to go by (of time)\textquotedblright{}.

\begin{quote}
\fa{اتوبوس} \fa{از} \fa{کنار} \fa{دریا} \fa{گذشت}.
\end{quote}

\begin{quote}
\emph{Otobus az kenâr-e daryâ gozasht.}
\end{quote}

\begin{quote}
The bus went past the sea.
\end{quote}

Present stem \emph{gozar-}: \emph{migozarad}. Of time: \emph{do sâl gozasht} — two years went by.

\begin{cousinweb}
\emph{gozashtan} is inherited: EtymDB traces it to a Middle Persian verb (written \emph{wtštn'}) meaning 'to pass, to cross, to die'.
\end{cousinweb}

\begin{grammarlens}[title={What you've built}]
One more word for telling, explaining and describing at B1.
\end{grammarlens}

\subsection*{Your turn}
\begin{itemize}
\raggedright
  \item \emph{Say it:} \emph{gozashtan}
  \item \emph{Say it:} \emph{gozashtan}, once more
  \item \emph{Say it:} a sentence of your own with \emph{gozashtan}
  \item \emph{Recall:} say the Persian for to follow, then the Persian for to approach, then say \emph{gozashtan} again
\end{itemize}

\begin{tcolorbox}[breakable,colback=teal!4,colframe=teal!35!black,title={Before you move on}]
What does \fa{گذشتن} mean? (\textquotedblleft{}To pass (by, through); to go by (of time)\textquotedblright{}.) Say it once more.
\end{tcolorbox}
\section[\texorpdfstring{\fa{کوچه} \fa{پس} \fa{کوچه} (kuche-pas-kuche)}{کوچه پس کوچه (kuche-pas-kuche)}]{\fa{کوچه} \fa{پس} \fa{کوچه} (kuche-pas-kuche) — the back streets, a maze of lanes}

\label{lesson:FA-C449-kuche-pas-kuche}



\begin{quote}
Before the new one: say the Persian for to approach, then the Persian for to pass.
\end{quote}

\subsection*{You'll want to know: \fa{کوچه} \fa{پس} \fa{کوچه}}
\textbf{\fa{کوچه} \fa{پس} \fa{کوچه}} (\emph{kuche-pas-kuche}) — \textquotedblleft{}the back streets, a maze of lanes\textquotedblright{}.

\begin{quote}
\fa{در} \fa{کوچه} \fa{پس} \fa{کوچه} \fa{های} \fa{شهر} \fa{قدیمی} \fa{قدم} \fa{زدیم}.
\end{quote}

\begin{quote}
\emph{Dar kuche-pas-kuche-hâ-ye shahr-e qadimi qadam zadim.}
\end{quote}

\begin{quote}
We walked through the back streets of the old town.
\end{quote}

Literally 'lane behind lane'. Usually in the plural.

\begin{cousinweb}
\emph{kuche-pas-kuche} is a native compound of \emph{kuche} 'lane', from \emph{ku} 'street', and \emph{pas} 'behind', inherited from Middle Persian \emph{pas} (EtymDB): lane after lane.
\end{cousinweb}

\begin{grammarlens}[title={What you've built}]
One more word for telling, explaining and describing at B1.
\end{grammarlens}

\subsection*{Your turn}
\begin{itemize}
\raggedright
  \item \emph{Say it:} \emph{kuche-pas-kuche}
  \item \emph{Say it:} \emph{kuche-pas-kuche}, once more
  \item \emph{Say it:} a sentence of your own with \emph{kuche-pas-kuche}
  \item \emph{Recall:} say the Persian for to approach, then the Persian for to pass, then say \emph{kuche-pas-kuche} again
\end{itemize}

\begin{tcolorbox}[breakable,colback=teal!4,colframe=teal!35!black,title={Before you move on}]
What does \fa{کوچه} \fa{پس} \fa{کوچه} mean? (\textquotedblleft{}The back streets, a maze of lanes\textquotedblright{}.) Say it once more.
\end{tcolorbox}
